%==============================================================================
\chapter{Results and Discussion}
\label{ch:results}
%==============================================================================

This chapter reports the experimental characterization of the sensor components
and the \nto\ detection results, followed by a discussion of the achieved
performance and its limitations.

\section{Tuning-Fork / ADM Characterization}
\label{sec:adm-char}

Report the measured resonance frequency $f_0$ and quality factor $Q$ of the QTF
(e.g.\ from a resonance sweep), and any dependence on gas pressure. These set the
operating modulation frequency and the integration bandwidth.

\section{Lock-In Amplifier Characterization}
\label{sec:lockin-char}

Report the lock-in characterization: response vs.\ time constant / filter
settings, noise floor, and the resulting detection bandwidth.

\section{QCL Characterization}
\label{sec:qcl-char}

Report the QCL output power and its wavelength tuning vs.\ temperature/current,
which maps drive settings onto the \nto\ absorption line.

\section{Optical Power Budget}
\label{sec:power-budget}

Present the measured transmission and power losses along the beam path and the
power delivered to the QTF, comparing with the design budget from
\cref{sec:alignment}.

\section{Nitrous Oxide Detection}
\label{sec:n2o-detection}

Show the photoacoustic signal as the laser is scanned across the selected \nto\
line, and the concentration calibration (signal vs.\ \nto\ concentration),
commenting on linearity.

\section{Sensitivity and Stability}
\label{sec:sensitivity}

Report the achieved sensitivity --- signal-to-noise ratio, normalized
noise-equivalent absorption (NNEA), and minimum detectable concentration --- and
the Allan-deviation analysis identifying the optimal integration time.

\section{Discussion}
\label{sec:discussion}

Interpret the results, compare with values reported in the literature, and
discuss limitations and sources of uncertainty (e.g.\ pulsed-mode spectral
broadening, alignment sensitivity, interferents).
